% Generated draft: requires mathematical and source review.
\documentclass[11pt]{article}
\usepackage[margin=25mm]{geometry}
\usepackage{fontspec}
\setmainfont{Latin Modern Roman}
\usepackage{amsmath,amssymb,mathtools,unicode-math}
\setmathfont{Latin Modern Math}
\usepackage{enumitem,needspace,xurl,hyperref}
\setlength{\parindent}{0pt}
\setlength{\parskip}{5pt}
\setlength{\emergencystretch}{3em}
\newcommand{\N}{\mathbb{N}}
\newcommand{\Z}{\mathbb{Z}}
\newcommand{\Q}{\mathbb{Q}}
\newcommand{\R}{\mathbb{R}}
\newcommand{\C}{\mathbb{C}}
\newcommand{\F}{\mathbb{F}}
\newcommand{\A}{\mathbb{A}}
\newcommand{\PP}{\mathbb{P}}
\newcommand{\E}{\mathbb{E}}
\newcommand{\eps}{\varepsilon}
\newcommand{\dd}{\,\mathrm{d}}
\begin{document}
\section*{UnsolvedMath ID: 6000007}
\textbf{Source problem code:} catalogue entry
\subsection{Definitions and mathematical statement}
A Riemannian metric \(g\) is locally conformally flat if near every point there are coordinates in which \(g=e^{2u}\sum_i(\dd x^i)^2\) for a smooth real function \(u\). A statistical manifold, in the usual Riemannian convention, is a triple \((M,g,\nabla)\) with \(\nabla\) torsion free and \(\nabla g\) totally symmetric. Its dual connection \(\nabla^*\) is determined by
\[
 Xg(Y,Z)=g(\nabla_XY,Z)+g(Y,\nabla_X^*Z).
\]

Kurose's problem asks to investigate the relationship between the geometry of a conformally flat Riemannian manifold and Matsuzoe's geometric divergence for a conformally-projectively flat statistical manifold. The excerpt does not define the particular conformal-projective equivalence, its flatness condition, or Matsuzoe's divergence. It also does not specify whether the Riemannian manifold is the underlying metric of that statistical manifold or is related to it by a separate construction.

For context, divergence-based statistical geometry often uses a two-point function \(D(x,y)\) whose diagonal mixed derivatives encode the metric and connection; a common convention is
\[
 g_{ij}(p)=-\left.\frac{\partial^2D}{\partial x^i\partial y^j}\right|_{x=y=p},\qquad
 \Gamma_{ijk}(p)=-\left.\frac{\partial^3D}{\partial x^i\partial x^j\partial y^k}\right|_{x=y=p}.
\]
This general convention does not identify Matsuzoe's specific geometric divergence, its normalization, positivity, or local/global domain, and is not substituted for the absent definition.

The request is a research theme rather than a stated equivalence, construction, invariant formula, or classification theorem. A precise target requires the original definitions and the intended relationship to be established. The supplied summary reports existing work on conformally flat statistical manifolds, Kurose divergence, Matsuzoe's constructions, and conformal-projective flatness, while saying no complete solution of this exact question was verified. That qualified open/triage assessment is retained without an independent current-status claim.
\subsection{Short English statement}
Clarify how conformally flat Riemannian geometry relates to Matsuzoe's divergence construction for conformally-projectively flat statistical manifolds. The source does not define the specific construction or a precise theorem to prove, so these remain explicit formulation gaps.
\subsection{Sources}
\begin{itemize}
\item Local registry: \url{data/unsolvedmath/problems.json}, permanent UnsolvedMath ID \(6000007\), source code catalogue entry; supplied statement, background, and research summary, updated 2026-08-06. Attributed to Takashi Kurose, 1998.
\item Furuhata, Matsuzoe, and Urakawa, \emph{Open Problems in Affine Differential Geometry and Related Topics} (1998), item 3(a), scanned PDF pages 1--2: \url{https://www.jstage.jst.go.jp/article/iis/4/2/4_2_125/_pdf/-char/en}. Supplied manual OCR from page images, accessed 2026-07-29.
\item The summary names Kurose divergence and Matsuzoe's geometric-divergence work without separate original-paper references. The supplied J-STAGE rights note identifies the article as CC BY 4.0.
\end{itemize}
\end{document}
